\section{Núcleo topológico de fase y generación del catálogo de semillas}
\label{act21:tpk:capitulo}

Las dos evaluaciones de APP adquieren eficacia generativa cuando se especifica
qué evaluación actúa, en qué posición se efectúa, cómo se desplaza esa posición
y qué información conserva cada transición. El \emph{Topological Phase Kernel},
abreviado TPK y denominado en español \emph{núcleo topológico de fase}, reúne
estas operaciones en una dinámica sobre estados enriquecidos. Su definición
comprende selección, transporte, actualización aritmética y conservación del
registro. La construcción finita desarrollada en este capítulo proporciona
las semillas sobre las que actuarán las elevaciones regionales y la prolongación
coinductiva.

El resultado principal de esta etapa es un catálogo generado y clasificado:
las condiciones iniciales de los dos cursores producen emisiones enteras;
su lectura ternaria determina palabras orientadas; la acción diedral organiza
esas palabras en órbitas. Cada paso conserva sus preimágenes y las relaciones
que permiten regresar a la descripción precedente. La cadena de cardinales
\[
 104\,976\longrightarrow468\longrightarrow243\subset729,
 \qquad 243\longrightarrow43,
\]
abrevia objetos diferentes. A continuación se construye cada objeto y se
justifica cada recuento, antes de utilizar cualquiera de ellos como dato de
la selección regional.

\subsection{Selección de fase, transporte y estado enriquecido}
\label{act21:tpk:operaciones}

\subsubsection{Selector trítico de fase operativa}

La fase nonádica toma sus valores en \(D_9=\{1,\ldots,9\}\). El selector
trítico es la función
\begin{equation}
 M_{\mathrm{ph}}:D_9\longrightarrow\{-1,0,+1\},\qquad
 M_{\mathrm{ph}}(g)=
 \begin{cases}
 +1,&g\in\{1,4,7\},\\
 -1,&g\in\{2,5,8\},\\
 0,&g\in\{3,6,9\}.
 \end{cases}
 \label{act21:tpk:selector}
\end{equation}
En la realización emisora, el valor \(+1\) activa la evaluación aditiva,
\(-1\) activa la evaluación multiplicativa y \(0\) registra el evento
neutral. La orientación cardinal del cursor es otra coordenada del estado.
Esta separación permite fijar simultáneamente qué operación se lee y en qué
dirección se transporta su punto de evaluación.

Para las transiciones numeradas por \(t\geq1\), la fase previa a la lectura es
\begin{equation}
 g_t=1+[(t-1)]_9,\qquad \epsilon_t=M_{\mathrm{ph}}(g_t),
 \label{act21:tpk:fase}
\end{equation}
donde \([n]_b\) designa el representante de \(n\) módulo \(b\) en
\(\{0,\ldots,b-1\}\). Una ventana de nueve transiciones contiene la sucesión
\((+1,-1,0,+1,-1,0,+1,-1,0)\).

\subsubsection{Cursores orientados y transporte cardinal}

Se utilizan índices \(I_9=\{0,\ldots,8\}\) para las posiciones y etiquetas
positivas \((a,b)=(i+1,j+1)\) para las evaluaciones APP. Un cursor finito es
\[
 c=(i,j;d)\in\mathcal S:=I_9^2\times\mathcal D,
 \qquad \mathcal D=\{N,E,S,O\}.
\]
Las cuatro direcciones actúan mediante
\[
 v_N=(-1,0),\quad v_E=(0,1),\quad
 v_S=(1,0),\quad v_O=(0,-1).
\]
La aplicación cardinal correspondiente a \(d\) es
\[
 T_d(i,j)=\bigl([i+(v_d)_1]_9,[j+(v_d)_2]_9\bigr).
\]
Así, \(M_{\mathrm{ph}}\) y \(T_d\) tienen dominios, codominios y acciones
diferentes. La activación trítica selecciona uno de los dos cursores;
el transporte cardinal modifica su posición en la dirección que ese cursor
ya conserva.

El levantamiento entero del cursor es
\(\widetilde c=(x,y;d)\in\mathbb Z^2\times\mathcal D\), con posición finita
\(([x]_9,[y]_9)\). Sus números de vuelta se recuperan en
\begin{equation}
 x=9\lfloor x/9\rfloor+[x]_9,\qquad
 y=9\lfloor y/9\rfloor+[y]_9.
 \label{act21:tpk:vueltas}
\end{equation}
Estas dos divisiones conciernen al transporte espacial. La división de una
evaluación aditiva o multiplicativa constituye un registro aritmético
distinto, que se conservará junto a ellas.

\subsubsection{Composición de la transición}

El estado de la realización finita incluye los dos cursores levantados,
sus direcciones, la fase, los acumuladores de la ventana activa, la lista
ordenada de bloques ya emitidos y el registro de transiciones. Cada registro
local contiene las posiciones previas, las evaluaciones enteras de ambas
hojas, sus residuos y cocientes, el régimen activo y las orientaciones.
Los campos de prefijo y frontera que aparezcan en una prolongación se
incorporarán al mismo estado con sus reglas de compatibilidad.

La transición se ordena como
\begin{equation}
 \mathcal U_t=\mathsf{Mem}_t\circ\mathsf{Tra}_t\circ\mathsf{Sel}_t.
 \label{act21:tpk:composicion}
\end{equation}
La selección determina el régimen y conserva la muestra previa a todo
desplazamiento. El transporte aplica el movimiento activo y la inversión
cardinal prescrita por el calendario. La actualización incorpora esa muestra
a los acumuladores y al registro; al terminar una ventana, emite su bloque.
La composición utiliza por tanto una muestra anterior al transporte aunque
su archivo definitivo se efectúe al final de la transición.

\subsection{Condiciones iniciales y enumeración reversible}
\label{act21:tpk:semillas}

Las hojas aditiva y multiplicativa disponen de cursores independientes. El
dominio de condiciones iniciales es el producto ordenado
\[
 \Omega_{\mathrm{seed}}=\mathcal S_+\times\mathcal S_\times,
 \qquad |\mathcal S|=9^2\cdot4=324.
\]
Por consiguiente,
\begin{equation}
 |\Omega_{\mathrm{seed}}|=324^2=104\,976.
 \label{act21:tpk:censo-inicial}
\end{equation}
El número \(81\) cuenta posiciones; \(324\) cuenta posiciones orientadas
de un cursor; \(104\,976\) cuenta parejas ordenadas de cursores. La elección
de una posición en APP conserva únicamente una parte de una condición
inicial del emisor.

Se fija la numeración
\[
 \operatorname{ind}(N)=0,\quad\operatorname{ind}(E)=1,\quad
 \operatorname{ind}(S)=2,\quad\operatorname{ind}(O)=3,
\]
y se define
\begin{equation}
 \nu(i,j;d)=4(9i+j)+\operatorname{ind}(d).
 \label{act21:tpk:indice-semilla}
\end{equation}

\begin{proposition}[Enumeración exacta de las condiciones iniciales]
\label{act21:tpk:enumeracion}
La aplicación \(\nu\) es una biyección de \(\mathcal S\) sobre
\(\{0,\ldots,323\}\). La pareja de índices se codifica de manera reversible
por \(324\nu_++\nu_\times\in\{0,\ldots,104\,975\}\).
\end{proposition}
\begin{proof}
La división por cuatro recupera
\[
 \operatorname{ind}(d)=[\nu]_4,\qquad
 j=[\lfloor\nu/4\rfloor]_9,\qquad i=\lfloor\nu/36\rfloor.
\]
Las cotas del dominio garantizan que estas tres cantidades pertenecen a
sus conjuntos de índices y reconstruyen \(\nu\) mediante
\eqref{act21:tpk:indice-semilla}. Para la pareja, la división euclídea por
\(324\) recupera cociente \(\nu_+\) y residuo \(\nu_\times\).
\end{proof}

La enumeración recorre el dominio completo antes de cualquier agrupación
regional. Los índices se utilizan para localizar las preimágenes de una
emisión, manteniendo explícita la condición que produjo cada trayectoria.

\subsection{Evaluación de las hojas y registro aritmético}
\label{act21:tpk:lecturas}

En una posición positiva \((a,b)\), las evaluaciones enteras son
\(A_+(a,b)=a+b\) y \(A_\times(a,b)=ab\). Para \(n>0\), se mantienen
\begin{equation}
 \rho_9(n)=1+[n-1]_9,\qquad
 q_9(n)=\lfloor(n-1)/9\rfloor,\qquad
 n=\rho_9(n)+9q_9(n).
 \label{act21:tpk:division-evaluacion}
\end{equation}
La activación aditiva incorpora \(\rho_9(A_+)\) al acumulador
correspondiente. La activación multiplicativa utiliza un observable del
residuo \(\rho_9(A_\times)\), definido a continuación.

\subsubsection{Exponente discreto en las unidades módulo nueve}

Las potencias de \(2\) módulo \(9\) recorren
\[
 2^0,2^1,\ldots,2^5\equiv1,2,4,8,7,5\pmod9,
 \qquad 2^6\equiv1\pmod9.
\]
El logaritmo discreto en este grupo de seis unidades se representa por
\begin{equation}
 \ell_9(1)=0,\quad\ell_9(2)=1,\quad\ell_9(4)=2,
 \quad\ell_9(8)=3,\quad\ell_9(7)=4,\quad\ell_9(5)=5.
 \label{act21:tpk:logaritmo}
\end{equation}
La ley emisora extiende el observable mediante
\(\ell_9(3)=\ell_9(6)=\ell_9(9)=0\). Para conservar su procedencia, la
muestra registrada incluye el residuo original y el indicador de pertenencia
al grupo de unidades. De ese modo, un valor observado cero puede proceder
de la unidad \(1\) o de alguno de los tres residuos no invertibles, y su
origen permanece recuperable.

\subsubsection{Orden de lectura y movimiento}

Las evaluaciones se efectúan en las posiciones anteriores a la transición.
Si \(\epsilon_t=+1\), se acumula la lectura aditiva y se desplaza su cursor.
Si \(\epsilon_t=-1\), se acumula el exponente discreto multiplicativo y se
desplaza el segundo cursor. Si \(\epsilon_t=0\), se incrementa el contador
neutral y se conservan ambas posiciones.

Al finalizar los pasos \(27\) y \(54\), las dos direcciones se sustituyen
por sus opuestas. Las ventanas se encadenan conservando posiciones y
orientaciones. Al inicio de cada ventana sólo se ponen a cero sus tres
acumuladores locales; el registro anterior permanece disponible.

\begin{proposition}[Inversión del transporte levantado]
\label{act21:tpk:inversion}
Sean \(a,b\in\{0,1\}\) los indicadores de activación e inversión de un
cursor. Si \(\operatorname{opp}\) intercambia \(N\leftrightarrow S\) y
\(E\leftrightarrow O\), la transformación
\[
 F_{a,b}(z,d)=\bigl(z+a v_d,\operatorname{opp}^{\,b}(d)\bigr)
\]
tiene inversa
\[
 F_{a,b}^{-1}(z',d')=
 \bigl(z'-a v_{\operatorname{opp}^{\,b}(d')},
       \operatorname{opp}^{\,b}(d')\bigr).
\]
La reducción espacial módulo nueve entrelaza esta transformación con el
transporte finito.
\end{proposition}
\begin{proof}
Aplicar dos veces la oposición recupera la dirección inicial. Una vez
recuperada esa dirección, la resta de \(a v_d\) deshace el desplazamiento.
Para la proyección espacial basta la igualdad
\([z+a v_d]_9=[[z]_9+a v_d]_9\), componente a componente. La inversión de
dirección conserva esta igualdad.
\end{proof}

La fase determina los indicadores utilizados por la inversa. El registro
espacial y el aritmético pueden así conservarse conjuntamente, incluso al
atravesar el borde del representante \(I_9^2\).

\subsection{Seis ventanas y emisión decimal por bloques}
\label{act21:tpk:emision}

\subsubsection{Acumulación en una ventana nonádica}

La ventana \(m\), con \(1\leq m\leq6\), comprende los pasos
\(9m-8,\ldots,9m\). Sean \(S_m\) la suma de sus residuos positivos
aditivos activos, \(E_m\) la suma de sus exponentes discretos activos y
\(Z_m\) su número de eventos neutrales.

\begin{lemma}[Multiplicidades de activación]
\label{act21:tpk:tres-eventos}
Cada ventana contiene tres activaciones aditivas, tres multiplicativas y
tres neutrales. En particular, \(Z_m=3\).
\end{lemma}
\begin{proof}
Las fases de una ventana recorren exactamente \(1,\ldots,9\). Cada uno
de los tres conjuntos de \eqref{act21:tpk:selector} tiene cardinal tres.
\end{proof}

Las seis ventanas forman \(54\) transiciones y producen seis bloques.
Esta longitud pertenece al emisor inicial aquí definido. La continuación
en profundidad actúa después sobre las palabras y sus registros compatibles.

\subsubsection{Regla de codificación decimal}

La ley emisora especifica las cifras
\begin{align}
 d_{m,1}&=[S_m]_{10},\notag\\
 d_{m,2}&=[S_m+E_m]_{10},\notag\\
 d_{m,3}&=[3S_m+5E_m+7Z_m]_{10},\notag\\
 U_m&=100d_{m,1}+10d_{m,2}+d_{m,3}.
 \label{act21:tpk:bloque-decimal}
\end{align}
Los pesos \(100,10,1\) son las posiciones centena, decena y unidad.
El subíndice \(10\) indica el residuo decimal, y cada bloque satisface
\(0\leq U_m\leq999\). La palabra emitida es la sucesión ordenada
\(\mathbf U=(U_1,\ldots,U_6)\), con ancho fijo de tres cifras por bloque.

Los coeficientes \(3,5,7\) constituyen los pesos explícitos de esta ley
emisora. Una vez fijada la regla, cada cifra se calcula a partir de los
acumuladores, y los resultados de este capítulo se deducen de esa regla.
El número \(1\) que aparece en su forma reducida tiene una procedencia
adicional concreta: \(7Z_m=21\equiv1\pmod{10}\).

Si \(s_m=[S_m]_{10}\) y \(e_m=[E_m]_{10}\), se obtiene
\begin{equation}
 G(s_m,e_m)=100s_m+10[s_m+e_m]_{10}
                   +[3s_m+5e_m+1]_{10}=U_m.
 \label{act21:tpk:acoplamiento-firmas}
\end{equation}
La letra \(e_m\) denota aquí una componente de la firma de exponentes
discretos. Su definición pertenece al emisor finito.

\begin{proposition}[Recuperación de las dos firmas]
\label{act21:tpk:recuperacion-firmas}
La aplicación \(G:\{0,\ldots,9\}^2\to\{0,\ldots,999\}\) es inyectiva.
Para un bloque \(U=100a+10b+c\) de su imagen, la inversa es
\[
 s=a,\qquad e=[b-a]_{10},\qquad
 c=[3s+5e+1]_{10}.
\]
En consecuencia, la palabra \(\mathbf U\) determina las dos firmas de
seis componentes \(\mathbf s\) y \(\mathbf e\).
\end{proposition}
\begin{proof}
Los dos últimos sumandos de \eqref{act21:tpk:acoplamiento-firmas} pertenecen a
\(\{0,\ldots,99\}\), de modo que la centena recupera \(s\). La decena
es \([s+e]_{10}\); al restar \(s\) y reducir módulo diez se recupera
\(e\). La unidad verifica la tercera congruencia. La reconstrucción se
aplica independientemente a las seis posiciones de la palabra.
\end{proof}

Los totales enteros \(S_m,E_m\) se recuperan de las lecturas registradas;
las centenas y decenas recuperan sus residuos. Por ejemplo, un total
\(S_m=15\) produce \(s_m=5\). Esta diferencia identifica exactamente
qué información conserva cada nivel de codificación.

\subsubsection{Codificación posicional de la palabra de seis bloques}

Los seis bloques admiten una codificación entera de ancho fijo,
\begin{equation}
 \mathscr N(\mathbf U)=\sum_{m=1}^{6}U_m1000^{6-m},
 \qquad0\leq\mathscr N(\mathbf U)<1000^6.
 \label{act21:tpk:entero-base1000}
\end{equation}
El primer bloque ocupa la posición de mayor peso y el sexto, la de unidades
de base mil. Un bloque \(058\) ocupa una posición completa con valor
entero \(58\); su ancho de tres cifras forma parte de la presentación
decimal de esa posición.

\begin{proposition}[Recuperación posicional de las emisiones]
\label{act21:tpk:inversa-base1000}
La codificación \eqref{act21:tpk:entero-base1000} determina la palabra completa
mediante
\[
 U_m=\left[\left\lfloor
 \frac{\mathscr N(\mathbf U)}{1000^{6-m}}\right\rfloor\right]_{1000},
 \qquad1\leq m\leq6.
\]
\end{proposition}
\begin{proof}
Al dividir por \(1000^{6-m}\), los bloques anteriores producen múltiplos
enteros de \(1000\), el bloque \(m\) produce \(U_m\), y la contribución
de los posteriores pertenece a \([0,1)\). La parte entera elimina esta
última contribución y el residuo módulo \(1000\) elimina las primeras.
\end{proof}

Esta codificación conserva las seis emisiones. Su dominio posee posiciones
ordenadas de base mil, mientras el dominio APP posee posiciones espaciales
módulo nueve. La composición mantiene ambos índices: \(m\) identifica una
ventana de emisión y \((i,j)\) identifica la celda visitada durante una
transición de esa ventana.

\subsection{Cálculo completo de una emisión que comienza por 501}
\label{act21:tpk:ejemplo501}

Se fijan las condiciones iniciales
\[
 c_{+,0}=(0,4;N),\qquad c_{\times,0}=(2,3;N),
 \qquad \nu_+=16,\quad\nu_\times=84.
\]
Las posiciones positivas iniciales son \((1,5)\) y \((3,4)\). El índice
cronológico previo es \(t=0\), la primera fase leída es \(g_1=1\) y el
registro está vacío. Estas coordenadas son las entradas
del cálculo; los bloques se obtienen ejecutando la recurrencia.

\begin{table}[htbp]
\caption{Primera ventana: lecturas previas y acumulación activa.}
\label{act21:tpk:tabla-501}
\small
\begin{tabular}{@{}rcccrrr@{}}
\toprule
\(t\)&\(\epsilon_t\)&posición \(+\)&posición \(\times\)&
\(\Delta S\)&\(\Delta E\)&\((S,E,Z)\)\\
\midrule
1&\(+1\)&\((1,5)\)&\((3,4)\)&6&0&\((6,0,0)\)\\
2&\(-1\)&\((9,5)\)&\((3,4)\)&0&0&\((6,0,0)\)\\
3&0&\((9,5)\)&\((2,4)\)&0&0&\((6,0,1)\)\\
4&\(+1\)&\((9,5)\)&\((2,4)\)&5&0&\((11,0,1)\)\\
5&\(-1\)&\((8,5)\)&\((2,4)\)&0&3&\((11,3,1)\)\\
6&0&\((8,5)\)&\((1,4)\)&0&0&\((11,3,2)\)\\
7&\(+1\)&\((8,5)\)&\((1,4)\)&4&0&\((15,3,2)\)\\
8&\(-1\)&\((7,5)\)&\((1,4)\)&0&2&\((15,5,2)\)\\
9&0&\((7,5)\)&\((9,4)\)&0&0&\((15,5,3)\)\\
\bottomrule
\end{tabular}
\notatabla{Las posiciones son etiquetas positivas de las celdas anteriores
al movimiento. Los eventos neutrales incrementan únicamente \(Z\).}
\end{table}

Las tres evaluaciones aditivas activas son
\[
 1+5=6,\qquad9+5=14=5+9,\qquad8+5=13=4+9,
\]
de modo que \(S_1=6+5+4=15\). Las tres evaluaciones multiplicativas activas
son \(3\cdot4=12\), \(2\cdot4=8\) y \(1\cdot4=4\). Sus residuos
positivos son \(3,8,4\), con observables \(0,3,2\); por tanto,
\(E_1=5\). Finalmente \(Z_1=3\). La regla decimal da
\[
 d_{1,1}=[15]_{10}=5,\quad
 d_{1,2}=[15+5]_{10}=0,\quad
 d_{1,3}=[45+25+21]_{10}=1,
\]
y se obtiene \(U_1=501\).

La segunda lectura aditiva anterior tiene cociente aritmético \(1\),
mientras su cursor levantado está en \((-1,4;N)\), con número de vuelta
vertical \(-1\). Son dos datos diferentes del mismo registro. Ambos
intervienen en la recuperación íntegra de la evaluación y de su posición.

\begin{table}[htbp]
\caption{Las seis ventanas de la condición inicial \((16,84)\).}
\label{act21:tpk:tabla-seis501}
\begin{tabular}{@{}rrrrrrrr@{}}
\toprule
\(m\)&pasos&\(S_m\)&\(E_m\)&\(Z_m\)&\(s_m\)&\(U_m\)&\([-U_m]_3\)\\
\midrule
1&1--9&15&5&3&5&501&0\\
2&10--18&6&5&3&6&614&1\\
3&19--27&24&5&3&4&498&0\\
4&28--36&21&5&3&1&169&2\\
5&37--45&12&5&3&2&272&1\\
6&46--54&12&5&3&2&272&1\\
\bottomrule
\end{tabular}
\end{table}

La palabra resultante y sus firmas son
\begin{align*}
 \mathbf U&=(501,614,498,169,272,272),\\
 \mathbf s&=(5,6,4,1,2,2),&
 \mathbf e&=(5,5,5,5,5,5).
\end{align*}
La inversión cardinal tras el paso \(27\) modifica los recorridos de las
tres ventanas finales. El bloque repetido \(272\) corresponde a dos
ventanas sucesivas con posiciones y registro cronológico propios.

Como segundo caso, la condición inicial mínima \((\nu_+,\nu_\times)=(0,0)\)
produce
\[
 \mathbf U=(252,109,252,903,850,850),\quad
 \mathbf s=(2,1,2,9,8,8),\quad
 \mathbf e=(3,9,3,1,7,7).
\]
En este caso, ambos cursores regresan a su posición y orientación iniciales
tras \(54\) transiciones. La fase nonádica completa seis vueltas, el
contador cronológico avanza \(54\) unidades y el registro contiene
\(54\) muestras. La lectura posicional del retorno y
la historia registrada describen propiedades distintas de la ejecución.

\subsection{Factorización del censo finito}
\label{act21:tpk:censo}

Las firmas \(f_+(\nu)=\mathbf s\) y \(f_\times(\nu)=\mathbf e\) se
calculan ejecutando el cursor correspondiente durante las seis ventanas.
La independencia de posiciones y direcciones entre los dos cursores da
\begin{equation}
 T_{\mathrm{em}}(\nu_+,\nu_\times)
   =G^{(6)}\bigl(f_+(\nu_+),f_\times(\nu_\times)\bigr),
 \label{act21:tpk:factorizacion}
\end{equation}
donde \(G^{(6)}\) aplica \eqref{act21:tpk:acoplamiento-firmas} en cada posición.
La factorización permite enumerar primero \(324+324\) trayectorias y
combinar después las firmas distintas, conservando la multiplicidad de
sus preimágenes.

\par\noindent\begin{minipage}{\linewidth}
\subsubsection{Imágenes de las dos firmas}

Las tablas siguientes incluyen todas las firmas. Su multiplicidad es el
número de condiciones iniciales que las producen; el índice mínimo localiza
una de esas condiciones. La preimagen completa se define por
\[
 f_\sigma^{-1}(a)=\{\nu\in\{0,\ldots,323\}:f_\sigma(\nu)=a\}.
\]
La recurrencia de las secciones anteriores determina cada elemento de este
conjunto mediante un cálculo entero finito.

\begin{table}[H]
\centering\fontsize{9.5}{11.4}\selectfont
\setlength{\abovecaptionskip}{8pt}\setlength{\belowcaptionskip}{6pt}
\renewcommand{\arraystretch}{1.12}
\caption{Las dieciocho firmas aditivas.}
\label{act21:tpk:firmas-aditivas}
\begin{tabular}{@{}lrr@{\qquad}lrr@{}}
\toprule
firma&multiplicidad&\(\nu_{\min}\)&firma&multiplicidad&\(\nu_{\min}\)\\
\midrule
\texttt{122564}&18&17&\texttt{564122}&18&16\\
\texttt{122898}&18&24&\texttt{645212}&18&4\\
\texttt{212645}&18&5&\texttt{654889}&18&33\\
\texttt{212988}&18&0&\texttt{889221}&18&13\\
\texttt{221456}&18&29&\texttt{889654}&18&32\\
\texttt{221889}&18&12&\texttt{898122}&18&25\\
\texttt{456221}&18&28&\texttt{898465}&18&20\\
\texttt{465898}&18&21&\texttt{988212}&18&1\\
\texttt{546988}&18&9&\texttt{988546}&18&8\\
\bottomrule
\end{tabular}

\par\vspace{12pt}

\caption{Las veintiséis firmas multiplicativas.}
\label{act21:tpk:firmas-multiplicativas}
\begin{tabular}{@{}lrr@{\qquad}lrr@{}}
\toprule
firma&multiplicidad&\(\nu_{\min}\)&firma&multiplicidad&\(\nu_{\min}\)\\
\midrule
\texttt{000000}&108&8&\texttt{555555}&24&84\\
\texttt{177393}&8&1&\texttt{555717}&8&14\\
\texttt{177555}&8&54&\texttt{555771}&8&18\\
\texttt{177771}&8&11&\texttt{555933}&8&24\\
\texttt{339555}&8&6&\texttt{717555}&8&12\\
\texttt{339771}&8&15&\texttt{717717}&8&21\\
\texttt{339933}&8&59&\texttt{717933}&8&19\\
\texttt{393177}&8&0&\texttt{771177}&8&9\\
\texttt{393393}&8&45&\texttt{771339}&8&13\\
\texttt{393555}&8&40&\texttt{771555}&8&16\\
\texttt{555177}&8&52&\texttt{933339}&8&57\\
\texttt{555339}&8&4&\texttt{933555}&8&26\\
\texttt{555393}&8&41&\texttt{933717}&8&17\\
\bottomrule
\end{tabular}
\end{table}
\end{minipage}\par


\begin{proposition}[Imagen y multiplicidades del emisor]
\label{act21:tpk:468}
La imagen \(\mathcal U_6=\operatorname{im}T_{\mathrm{em}}\) contiene
\(468\) palabras distintas. Sus preimágenes se distribuyen en \(432\)
fibras de cardinal \(144\), \(18\) de cardinal \(432\) y \(18\) de
cardinal \(1944\).
\end{proposition}
\begin{proof}
La evaluación exhaustiva del algoritmo de la sección~\ref{act21:tpk:algoritmo}
sobre los índices \(0,\ldots,323\) da las dos tablas anteriores.
La primera partición tiene \(18\) clases de cardinal \(18\), y la segunda
\(24\) clases de cardinal \(8\), una de cardinal \(24\) y una de
cardinal \(108\). Las sumas \(18\cdot18=324\) y
\(24\cdot8+24+108=324\) comprueban la masa total de ambas particiones.
La proposición~\ref{act21:tpk:recuperacion-firmas} hace inyectivo el acoplamiento
de sus imágenes; por tanto, hay \(18\cdot26=468\) emisiones.
Por \eqref{act21:tpk:factorizacion}, la preimagen de una pareja de firmas es
el producto de sus dos preimágenes. Se obtienen \(18\cdot24=432\)
productos de cardinal \(18\cdot8=144\), \(18\) de cardinal
\(18\cdot24=432\) y \(18\) de cardinal \(18\cdot108=1944\).
Finalmente,
\[
 432\cdot144+18\cdot432+18\cdot1944=104\,976.
\]
La parte enumerativa es una comprobación finita exacta de la recurrencia;
la factorización explica por qué ese recuento cubre todas las parejas.
\end{proof}

La emisión del ejemplo \ref{act21:tpk:ejemplo501} tiene una fibra de cardinal
\(18\cdot24=432\). El ejemplo de índices \((0,0)\) tiene una fibra de
cardinal \(18\cdot8=144\). En ambos casos la palabra emitida recupera
las firmas, mientras la identificación de una condición inicial particular
utiliza su índice o su registro de trayectoria.

\subsection{Lectura ternaria orientada y conservación de las fibras}
\label{act21:tpk:reduccion-ternaria}

La reducción orientada actúa separadamente sobre cada bloque:
\begin{equation}
 \Psi:\mathcal U_6\longrightarrow\mathbb F_3^6,\qquad
 \Psi(U_1,\ldots,U_6)=([-U_1]_3,\ldots,[-U_6]_3).
 \label{act21:tpk:psi}
\end{equation}
Su signo fija una orientación del catálogo. La identificación balanceada
\(0\mapsto0\), \(1\mapsto+1\), \(2\mapsto-1\) se aplica después,
componente a componente.

El módulo nueve organiza las evaluaciones y posiciones de APP; los residuos
módulo diez forman cada cifra; la base mil reúne tres posiciones decimales;
el módulo tres determina la firma ternaria de cada bloque. Estos cuatro usos
tienen objetos y funciones definidos. En particular, \(\Psi\) lee seis
bloques y produce seis trits; una conversión posicional del entero formado
por dieciocho cifras sería otra operación.

Para las dos condiciones calculadas se obtiene
\begin{align*}
 \Psi(501,614,498,169,272,272)&=(0,1,0,2,1,1),\\
 \Psi(252,109,252,903,850,850)&=(0,2,0,0,2,2).
\end{align*}
El emisor decimal está definido antes de esta reducción. La función de
\(\Psi\) es producir el objeto ternario sobre el que actúan la simetría
orbital y los selectores regionales. Esa función explica su posición en la
genealogía: enlaza una emisión ya calculada con una clasificación estructural.

La imagen \(\mathcal W_6:=\Psi(\mathcal U_6)\) contiene exactamente
\(243\) palabras dentro del conjunto ambiente de \(3^6=729\) palabras.
La enumeración del algoritmo final da el histograma
\begin{equation}
\begin{array}{c|rrrrrr}
 |\Psi^{-1}(w)|&1&2&3&4&5&6\\\hline
 \text{número de palabras }w&132&66&18&3&6&18.
\end{array}
\label{act21:tpk:fibras-ternarias}
\end{equation}
Las sumas de ambas lecturas son
\[
132+66+18+3+6+18=243,
\qquad132+2\cdot66+3\cdot18+4\cdot3+5\cdot6+6\cdot18=468.
\]
La primera cuenta palabras ternarias; la segunda recupera el número de
emisiones distribuidas entre sus fibras.

\begin{example}[Dos emisiones con la misma firma ternaria]
\label{act21:tpk:colision}
Las emisiones
\[
 (498,501,614,252,252,109),\qquad
 (870,810,923,252,252,109)
\]
son distintas y ambas tienen imagen \((0,0,1,0,0,2)\) por \(\Psi\).
Cada una tiene \(144\) condiciones iniciales. La palabra ternaria conserva
una clase de lectura; las emisiones y sus preimágenes distinguen los datos
que confluyen en esa clase.
\end{example}

La conservación de fibras se expresa concretamente mediante
\[
 \mathcal F_w=\{\mathbf U\in\mathcal U_6:\Psi(\mathbf U)=w\},
 \qquad
 \Omega_{\mathbf U}=T_{\mathrm{em}}^{-1}(\mathbf U).
\]
Por tanto, la preimagen completa de \(w\) en el dominio inicial es la
unión disjunta de los conjuntos \(\Omega_{\mathbf U}\), con
\(\mathbf U\in\mathcal F_w\). El registro enriquecido puede conservar la
emisión, sus firmas y la condición inicial junto a su lectura ternaria.
La inversa exacta se refiere entonces al dato registrado con su tipo;
la palabra de seis trits aislada mantiene el alcance preciso de
\eqref{act21:tpk:psi}.

\subsection{Acción diedral y las cuarenta y tres órbitas}
\label{act21:tpk:orbitas}

Sobre seis coordenadas se definen las permutaciones
\begin{align}
 r(u_1,u_2,u_3,u_4,u_5,u_6)
   &=(u_3,u_1,u_2,u_5,u_6,u_4),\\
 s(u_1,u_2,u_3,u_4,u_5,u_6)
   &=(u_4,u_5,u_6,u_1,u_2,u_3).
 \label{act21:tpk:generadores-diedrales}
\end{align}
La primera rota los dos grupos de tres coordenadas con sentidos conjugados;
la segunda intercambia ambos grupos. Su composición satisface
\(r^3=s^2=1\) y \(srs=r^{-1}\). Generan así una acción del grupo diedral
\(D_3\), de seis elementos.

\begin{proposition}[Equivariancia y estabilidad del catálogo]
\label{act21:tpk:equivarianza}
La reducción \(\Psi\) conmuta con la acción de \(D_3\). Los conjuntos
\(\mathcal U_6\) y \(\mathcal W_6\) son estables bajo sus generadores.
\end{proposition}
\begin{proof}
La igualdad \(\Psi(g\mathbf U)=g\Psi(\mathbf U)\) se verifica en cada
coordenada porque \(g\) permuta posiciones y \(\Psi\) aplica la misma
reducción a cada posición. La estabilidad del catálogo constituye otra
comprobación: al generar sus \(468\) elementos mediante
\eqref{act21:tpk:factorizacion}, se verifica que \(r\mathbf U\) y
\(s\mathbf U\) pertenecen de nuevo a ese conjunto. El algoritmo de la
sección~\ref{act21:tpk:algoritmo} efectúa estas comprobaciones de pertenencia.
La equivariancia transmite después la estabilidad a su imagen ternaria.
\end{proof}

\begin{proposition}[Descomposición orbital completa]
\label{act21:tpk:43}
El conjunto \(\mathcal W_6\) se descompone en \(43\) órbitas: \(38\)
de cardinal seis y \(5\) de cardinal tres.
\end{proposition}
\begin{proof}
Para cada \(w\in\mathcal W_6\) se forma
\[
 \mathcal O(w)=\{w,rw,r^2w,sw,srw,sr^2w\}.
\]
Se escoge como representante la palabra lexicográficamente mínima de ese
conjunto. Dos palabras tienen el mismo representante exactamente cuando
pertenecen a la misma órbita. El algoritmo genera la lista de representantes
de la tabla~\ref{act21:tpk:representantes}; sus cardinales suman
\(38\cdot6+5\cdot3=243\). Como cada palabra del catálogo participa en
esta operación y el catálogo es estable, la lista cubre toda la imagen.
\end{proof}

\begin{table}[htbp]
\caption{Representantes mínimos y cardinales de las cuarenta y tres órbitas.}
\label{act21:tpk:representantes}
\small
\begin{tabular}{@{}lr@{\qquad}lr@{\qquad}lr@{}}
\toprule
representante&cardinal&representante&cardinal&representante&cardinal\\
\midrule
\texttt{001002}&6&\texttt{002112}&6&\texttt{012222}&6\\
\texttt{001011}&6&\texttt{002211}&6&\texttt{021022}&6\\
\texttt{001012}&6&\texttt{002220}&6&\texttt{021121}&6\\
\texttt{001021}&6&\texttt{011021}&6&\texttt{021202}&6\\
\texttt{001022}&6&\texttt{011112}&6&\texttt{022022}&3\\
\texttt{001100}&3&\texttt{011120}&6&\texttt{022112}&6\\
\texttt{001101}&6&\texttt{011201}&6&\texttt{022121}&6\\
\texttt{001110}&6&\texttt{011211}&6&\texttt{022202}&3\\
\texttt{001112}&6&\texttt{011220}&6&\texttt{022211}&6\\
\texttt{001202}&6&\texttt{011222}&6&\texttt{022222}&6\\
\texttt{001210}&6&\texttt{012112}&6&\texttt{112112}&3\\
\texttt{001211}&6&\texttt{012121}&6&\texttt{112211}&3\\
\texttt{001220}&6&\texttt{012202}&6&\texttt{112222}&6\\
\texttt{001222}&6&\texttt{012210}&6&&\\
\texttt{002012}&6&\texttt{012220}&6&&\\
\bottomrule
\end{tabular}
\end{table}

Las órbitas heredan invariantes enteros de sus fibras. Para cada palabra,
\(n(w)=|\mathcal F_w|\) cuenta emisiones, mientras
\[
 M(w)=\sum_{\mathbf U\in\mathcal F_w}|\Omega_{\mathbf U}|
\]
cuenta condiciones iniciales. El censo de símbolos
\(c(w)=(n_0,n_1,n_2)\) registra cuántas veces aparece cada clase ternaria.
Estos datos, junto con la orientación y el soporte de los recorridos,
proporcionan los predicados de la selección regional posterior. La
clasificación conserva así más información que la forma de una palabra
representante.

\subsection{Algoritmo autocontenido de enumeración y control}
\label{act21:tpk:algoritmo}

La enumeración utiliza enteros, listas finitas y las operaciones ya definidas.
El siguiente procedimiento precisa su ejecución; \(\sigma\) toma los valores
\(+\) y \(\times\).

\begin{enumerate}
\item Para cada \(\nu=0,\ldots,323\), recuperar \((i,j;d)\) mediante
la inversa de \eqref{act21:tpk:indice-semilla}. Iniciar una lista vacía de firmas.
\item Para cada ventana \(m=1,\ldots,6\), iniciar el acumulador \(A=0\).
Recorrer \(t=9m-8,\ldots,9m\), calculando \(\epsilon_t\) por
\eqref{act21:tpk:fase}.
\item Si \(\sigma=+\) y \(\epsilon_t=+1\), añadir
\(\rho_9(i+j+2)\) a \(A\) y actualizar
\((i,j)\leftarrow T_d(i,j)\). Si \(\sigma=\times\) y
\(\epsilon_t=-1\), añadir
\(\ell_9(\rho_9((i+1)(j+1)))\) y efectuar el mismo transporte.
\item Al terminar \(t=27\) o \(t=54\), sustituir \(d\) por su opuesta.
Al concluir la ventana, añadir \([A]_{10}\) a la lista de firmas.
Conservar la posición y la dirección para la siguiente ventana.
\item Agrupar los \(324\) índices por igualdad de sus listas de seis
componentes. Estas dos agrupaciones dan \(f_+\), \(f_\times\) y todas
sus preimágenes.
\item Para cada pareja de firmas distintas, aplicar \(G^{(6)}\).
Asignar a la emisión la multiplicidad producto de las dos preimágenes.
Comprobar la recuperación de firmas en cada bloque y la suma de
multiplicidades \(104\,976\).
\item Aplicar \(\Psi\) a cada emisión y agrupar por igualdad de palabras.
Contar los \(243\) valores y las multiplicidades de
\eqref{act21:tpk:fibras-ternarias}.
\item Aplicar \(r,s\) a cada emisión y verificar pertenencia al catálogo.
Para cada palabra ternaria, formar sus seis imágenes diedrales, eliminar
repeticiones y registrar su mínimo lexicográfico. Contar las \(43\)
órbitas y los dos tamaños de la tabla~\ref{act21:tpk:representantes}.
\end{enumerate}

La ejecución termina: el primer tramo evalúa \(648\) trayectorias de
\(54\) transiciones; el segundo combina \(468\) parejas de firmas; el
tercero agrupa un conjunto de \(243\) palabras. Una comprobación directa
adicional puede ejecutar las \(104\,976\) parejas y contrastar la
factorización \eqref{act21:tpk:factorizacion}; su alcance coincide con el mismo
dominio finito.

\subsection{Conservación del registro y función de las semillas}
\label{act21:tpk:conclusion}

Sea \(\mathsf{Sample}(q_t)\) la muestra previa a la transición y
\(\mathcal M_t\) el registro acumulado. La actualización de memoria es
\begin{equation}
 \mathcal M_{t+1}=\mathcal M_t\mathbin{\Vert}\mathsf{Sample}(q_t),
 \label{act21:tpk:archivo}
\end{equation}
donde \(\Vert\) denota concatenación ordenada. La inducción da
\[
 \mathcal M_n=\mathcal M_0\mathbin{\Vert}
 (\mathsf{Sample}(q_0),\ldots,\mathsf{Sample}(q_{n-1})).
\]
Los acumuladores se reconstruyen sumando las contribuciones activas de
cada segmento de nueve registros. Sus cocientes, las vueltas de los
cursores y las orientaciones permanecen asociados a las mismas transiciones.
La ejecución determina ese registro mediante la recurrencia; el archivo
preserva la procedencia concreta de las lecturas.

Al finalizar esta etapa están construidos el dominio completo de condiciones
iniciales, la emisión decimal, sus dos firmas recuperables, la imagen ternaria,
sus fibras y su clasificación orbital. Los censos \(104\,976\), \(468\),
\(243\), \(729\) y \(43\) tienen ahora dominio y mecanismo de obtención.
Las selecciones regionales recibirán estas estructuras con sus invariantes,
y las elevaciones actuarán sobre palabras orientadas con procedencia
conservada. Esa continuidad convierte el catálogo finito en antecedente
constructivo del desarrollo coinductivo.
